\problem{CP01}{跨越平行线的随机针}{几何概率与交点数}
平面上画间距为 $h>0$ 的平行直线。投下一根长度为 $l>0$ 的针，假设针中心到最近直线的距离在 $[0,h/2]$ 上均匀，针与平行线的锐角在 $[0,\pi/2]$ 上均匀，且二者独立。求针至少碰到一条线的概率和针与直线的交点数期望；针长可以超过线距。
\solution 记角度为 $\Theta$、中心距离为 $D$、长度比为 $r=l/h$。针的法向投影长度为 $l\sin\Theta$，碰线等价于 $D\le l\sin\Theta/2$。条件化于角度，
\[
p(r)=\frac2\pi\int_0^{\pi/2}\min(1,r\sin\theta)\dd\theta.
\]
若 $r\le1$，积分直接给出 $p(r)=2r/\pi$。若 $r>1$，以 $\theta_0=\arcsin(1/r)$ 分开积分，得到
\[
\boxed{p(r)=\begin{cases}
2r/\pi,&0<r\le1,\\
\dfrac2\pi\left(r-\sqrt{r^2-1}+\arccos(1/r)\right),&r>1.
\end{cases}}
\]
再求交点数 $J$。给定角度，针投影为长度 $b=l\sin\theta$ 的区间，直线位置构成间距 $h$ 的格点。交点数在中心位置关于最近直线反射时不变，所以距离均匀的模型可等价地用中心在一个周期内均匀平移来计算。此时，任意长度 $b$ 的区间所覆盖格点数的均值为 $b/h$：把各格点的被覆盖指示函数积分并相加，恰好覆盖长度为 $b$ 的平移范围。故
\[
\boxed{\E J=\E\left(\frac{l\sin\Theta}{h}\right)=\frac{2l}{\pi h}.}
\]
当 $l\le h$ 时，几乎必然至多一个交点，故概率与期望相等；长针可能跨多条线，两者必须区分。
\insight “随机投针”必须指定位置和方向的分布。先条件化角度，再求几何长度比例，可以同时处理短针和长针，而不依赖图形直觉。

\problem{CP02}{避免单色集合的二色着色}{概率方法与条件期望}
有限点集 $V$ 上给定 $m$ 个子集，每个子集至少含 $k\ge2$ 个点。若 $m<2^{k-1}$，证明存在一种红蓝二色着色，使每个指定子集都同时含两种颜色。进一步说明如何逐点确定颜色，而不必穷举全部着色。
\solution 先独立等概率把各点染成红或蓝。含 $s$ 个点的子集全红或全蓝的概率为 $2^{1-s}\le2^{1-k}$。令 $Z$ 为单色指定子集的个数，则
\[
\E Z\le m2^{1-k}<1.
\]
若每种着色都有单色子集，必有 $Z\ge1$，与期望小于 $1$ 矛盾。因此存在 $Z=0$ 的着色。联合界还给出成功概率至少为 $1-m2^{1-k}>0$。

可把这一存在性证明变成逐点选择。每次已有部分颜色固定时，计算当前条件期望。下一个点随机取红、蓝的两种条件期望的平均，等于当前值，因此至少一种颜色不会使条件期望增加，固定为该颜色即可。

这些条件期望可以直接计算。一个子集若已出现两色，单色概率为零；若已固定点均同色、尚有 $u$ 个未染点，单色概率为 $2^{-u}$；若尚无点固定，概率仍为 $2^{1-s}$。将各子集概率相加即可。所有点固定后，条件期望就是整数 $Z$，仍小于 $1$，故必为零。
\insight 随机性不仅用于计算概率，也能证明确定对象存在。每一步保留较小的条件期望，就把随机证明转成了确定选择规则；这里给出的数量条件是充分条件，不必是必要条件。

\problem{CP03}{一次一密与密钥复用}{条件独立与信息泄漏}
消息与密钥取值于 $\mathbb Z_q=\{0,\ldots,q-1\}$，$q\ge2$，加密为 $C=M+K\pmod q$。密钥 $K$ 与消息 $M$ 独立。证明均匀密钥使 $C$ 与 $M$ 独立；若消息每个值均有正概率，证明完全保密也要求密钥均匀。再令两条消息独立均匀，复用同一个独立均匀密钥得到 $C_1,C_2$，求给定两份密文后消息对的条件分布，并指出泄漏的确定信息。
\solution 若密钥均匀，对任意 $m,c$，
\[
\Prob(C=c\mid M=m)=\Prob(K=c-m)=\frac1q.
\]
条件分布不依赖消息，所以密文均匀且独立于消息，消息先验不必均匀。

反过来，若密文独立于消息且每个消息值均有正概率，则对固定 $c$，$\Prob(K=c-m)=\Prob(C=c)$ 对所有 $m$ 相等。随着 $m$ 遍历 $\mathbb Z_q$，$c-m$ 遍历所有密钥值，因而各密钥概率相等，只能都为 $1/q$。满支持假设不可省略，例如恒定消息与任何密文都独立。

复用时，对给定密文 $(c_1,c_2)$，每个密钥 $k$ 唯一对应消息对 $(c_1-k,c_2-k)$。原先每个三元组 $(M_1,M_2,K)$ 概率为 $q^{-3}$，所以该密文概率为 $q\cdot q^{-3}=q^{-2}$。于是
\[
\boxed{\begin{aligned}
&\Prob(M_1=m_1,M_2=m_2\mid C_1=c_1,C_2=c_2)\\
&\qquad=\begin{cases}
1/q,&m_2-m_1=c_2-c_1\pmod q,\\
0,&\text{其他情形}.
\end{cases}\end{aligned}}
\]
单条消息的后验仍均匀，但差值 $M_2-M_1$ 已被密文差值完全暴露。消息对从 $q^2$ 个可能缩减为 $q$ 个，不能再称为联合完全保密。
\insight 单份密文的保密性不保证多份密文的联合保密性。重复使用密钥后，相减消去密钥，正是依赖结构泄漏信息的原因。
